\chapter{The Schwartz Space}
\label{ch:iii14}

The Schwartz space packages smoothness and rapid decay into a Fréchet-type test-function space stable under Fourier transformation. It is the natural common domain for classical transform calculus and tempered distributions.

\section*{Learning goals}
After this chapter, the reader should be able to:
\begin{itemize}
\item verify that a function belongs to the Schwartz space by bounding all polynomially weighted derivatives;
\item show that differentiation, multiplication by polynomials, and Fourier transformation preserve the Schwartz space;
\item construct seminorm estimates controlling standard Schwartz-space operations;
\item distinguish rapid decay from compact support and from mere integrability;
\item use density of Schwartz functions in relevant \(L^p\) or \(L^2\) settings where established;
\item produce examples that lie in one standard function space but not another by comparing decay and smoothness requirements;
\end{itemize}
\section*{Conceptual roadmap}
\[
\boxed{\text{structure}\;\longrightarrow\;\text{estimate}\;\longrightarrow\;\text{limit or representation}\;\longrightarrow\;\text{application}.}
\]

\paragraph{Mathematical setting.}
Here $\mathcal S(\mathbb R^n)$ consists of smooth functions with every seminorm $p_{\alpha,\beta}(f)=\sup_x|x^\alpha D^\beta f(x)|$ finite, for multi-indices $\alpha,\beta$. Its topology is generated by this countable family. Stability estimates bound each output seminorm by a finite sum of input seminorms; a seminorm itself is not a linear combination. Density in $L^p(\mathbb R^n)$ is meant for $1\le p<\infty$.

\section{Rapid decay}
A Schwartz function decays faster than every polynomial together with all derivatives.

\section{Schwartz seminorms}
Quantities $\sup_x|x^\alpha D^\beta f(x)|$ encode the topology.


% BEGIN VOL03-EXPANSION III14-example-01
\begin{example}[Gaussian is Schwartz]\label{ex:iii14-expand-01}
Every derivative of \(e^{-x^2}\) is a polynomial times \(e^{-x^2}\).
Therefore for every pair of nonnegative integers \(m,k\),
\[
\sup_x |x^m f^{(k)}(x)|<\infty,
\]
because exponential decay dominates polynomial growth. This verifies all
Schwartz seminorms directly.
\end{example}
% END VOL03-EXPANSION III14-example-01
\section{Algebraic stability}
Multiplication by polynomials and differentiation preserve Schwartz space.

\section{Translation and modulation}
Translations and frequency modulations preserve the class.

\section{Fourier invariance}
The Fourier transform maps Schwartz space bijectively to itself.


% BEGIN VOL03-EXPANSION III14-example-02
\begin{example}[Fourier invariance through identities]\label{ex:iii14-expand-02}
Frequency differentiation corresponds to multiplication by powers of \(x\),
while frequency multiplication corresponds to spatial differentiation. Thus
every weighted derivative of \(\widehat f\) is the transform of an integrable
Schwartz expression built from \(f\). The \(L^1\) transform bound then controls
all Schwartz seminorms of \(\widehat f\).
\end{example}
% END VOL03-EXPANSION III14-example-02
\section{Convolution}
Schwartz functions convolve to Schwartz functions. Convolution of a Schwartz function with an arbitrary L1 function is smooth but need not have rapid decay; all polynomial absolute moments of that L1 function are a sufficient additional condition for Schwartz decay.

\section{Density}
Schwartz space is dense in $L^p$ for finite $p$ and in many Sobolev spaces.

\section{Role in distribution theory}
Its rapid decay makes pairing with polynomial-growth objects well-defined, leading to tempered distributions.


% BEGIN VOL03-EXPANSION III14-example-03
\begin{example}[Why tempered distributions use Schwartz tests]\label{ex:iii14-expand-03}
A polynomial-growth function can be paired with a Schwartz test function because
rapid decay beats every polynomial. For example
\[
T(\phi)=\int(1+x^2)^N\phi(x)\,dx
\]
is finite and controlled by Schwartz seminorms. This lets distribution theory
include slowly growing objects while preserving Fourier transform operations.
\end{example}
% END VOL03-EXPANSION III14-example-03
\section{Core structural results}

\begin{theorem}[Fourier invariance]\label{thm:iii14-01}
If $f\in\mathcal S$, then $\widehat f\in\mathcal S$.
\begin{proof}
Differentiation in frequency corresponds to multiplication by powers of $x$, while multiplication by powers of $\xi$ corresponds to derivatives of $f$; all such weighted derivatives are bounded and integrable for Schwartz functions.
\end{proof}
\end{theorem}

\begin{theorem}[Algebraic stability]\label{thm:iii14-02}
If $f\in\mathcal S$, then $x^\alpha D^\beta f\in\mathcal S$.
\begin{proof}
By the Leibniz rule, every Schwartz seminorm of this expression is bounded by a finite sum of higher Schwartz seminorms of $f$ with nonnegative constant coefficients.
\end{proof}
\end{theorem}

\begin{theorem}[Schwartz convolution closure]\label{thm:iii14-03}
If $f,g\in\mathcal S$, then $f*g\in\mathcal S$.
\begin{proof}
Differentiate under the integral and distribute polynomial weights between $x-y$ and $y$; rapid decay of both factors controls every resulting seminorm.
\end{proof}
\end{theorem}

\section{Worked examples}

\begin{example}[Polynomial times Gaussian is Schwartz]\label{ex:iii14-worked-01}
Let \(f(x)=x^me^{-\pi x^2}\). Every derivative has the form
\(P(x)e^{-\pi x^2}\) for a polynomial \(P\). Multiplying by any power \(x^k\)
still leaves a polynomial times a Gaussian, which is bounded. Hence
\(f\in\mathcal S(\mathbb R)\).
\end{example}

\begin{example}[A Schwartz seminorm calculation]\label{ex:iii14-worked-02}
For
\[
p_{2,1}(f)=\sup_x(1+|x|)^2|f'(x)|,
\]
the Gaussian derivative \(f'(x)=-2\pi xe^{-\pi x^2}\) makes
\[
p_{2,1}(f)=\sup_x 2\pi(1+|x|)^2|x|e^{-\pi x^2}<\infty.
\]
The exact maximum is irrelevant; finiteness is the topology-defining fact.
\end{example}

\begin{example}[Fourier transform preserves the Schwartz class]\label{ex:iii14-worked-03}
Differentiation of \(\widehat f\) corresponds to moments of \(f\), while
multiplication by \(\xi\) corresponds to transforms of derivatives. Because
all derivatives and polynomial moments of a Schwartz function are integrable
and rapidly controlled, every Schwartz seminorm of \(\widehat f\) is finite.
\end{example}

\begin{example}[Compact support implies Schwartz after smoothness]\label{ex:iii14-worked-04}
If \(\varphi\in C_c^\infty(\mathbb R^n)\), then every derivative is
bounded and supported in one compact set. Therefore
\[
\sup_x |x^\alpha D^\beta\varphi(x)|<\infty
\]
for every \(\alpha,\beta\), so \(C_c^\infty\subset\mathcal S\).
\end{example}

\section{Solved dossiers}
Each dossier combines several ideas from the chapter in a longer argument. The aim is to make hypotheses, calculations, and proof strategy explicit.

\begin{problem}[Polynomial times Gaussian is Schwartz]\label{prob:iii14-dossier-01}
Prove \(x^me^{-\pi x^2}\in\mathcal S(\mathbb R)\).
\end{problem}
\begin{solution}
Repeated differentiation produces polynomial factors only, and Gaussian
decay dominates every polynomial after multiplication by any additional
power of \(x\).
\end{solution}

\begin{problem}[A Schwartz seminorm calculation]\label{prob:iii14-dossier-02}
Show a representative Schwartz seminorm of \(e^{-\pi x^2}\) is finite.
\end{problem}
\begin{solution}
Differentiate the Gaussian. The result is a polynomial times the same
Gaussian, whose supremum after any polynomial weight is finite.
\end{solution}

\begin{problem}[Fourier transform preserves the Schwartz class]\label{prob:iii14-dossier-03}
Explain structurally why \(\mathcal F:\mathcal S\to\mathcal S\).
\end{problem}
\begin{solution}
Frequency derivatives are transforms of polynomial moments, and frequency
weights are transforms of spatial derivatives. Schwartz control supplies both
families simultaneously.
\end{solution}

\begin{problem}[Compact support implies Schwartz after smoothness]\label{prob:iii14-dossier-04}
Prove every smooth compactly supported function is Schwartz.
\end{problem}
\begin{solution}
All derivatives vanish outside a compact set, and polynomial weights are
bounded on that compact set.
\end{solution}

\begin{problem}[Seminorm family]\label{prob:iii14-dossier-05}
State a standard family of seminorms defining the Schwartz topology.
\end{problem}
\begin{solution}
For example
\(p_{\alpha,\beta}(f)=\sup_x|x^\alpha D^\beta f(x)|\) for all multi-indices
\(\alpha,\beta\).
\end{solution}

\begin{problem}[Schwartz implies \(L^p\)]\label{prob:iii14-dossier-06}
Why is a Schwartz function in \(L^p\) for every \(1\le p\le\infty\)?
\end{problem}
\begin{solution}
Choose polynomial decay faster than \(n/p\); the Schwartz bounds dominate
the function by an integrable power tail.
\end{solution}

\begin{problem}[Derivative stability]\label{prob:iii14-dossier-07}
Show \(D^\gamma f\in\mathcal S\) when \(f\in\mathcal S\).
\end{problem}
\begin{solution}
The defining seminorms of \(D^\gamma f\) are among the higher-derivative
seminorms already finite for \(f\).
\end{solution}

\begin{problem}[Multiplication by polynomials]\label{prob:iii14-dossier-08}
Show \(x^\gamma f\in\mathcal S\) when \(f\in\mathcal S\).
\end{problem}
\begin{solution}
Leibniz differentiation produces finitely many polynomial-weighted
derivatives of \(f\), all controlled by Schwartz seminorms.
\end{solution}

\begin{problem}[Fourier invariance proof dossier]\label{prob:iii14-dossier-09}
Prove the following structural statement and identify the step where its main hypothesis is used: If $f\in\mathcal S$, then $\widehat f\in\mathcal S$.
\end{problem}
\begin{solution}
Differentiation in frequency corresponds to multiplication by powers of $x$, while multiplication by powers of $\xi$ corresponds to derivatives of $f$; all such weighted derivatives are bounded and integrable for Schwartz functions. This proof also explains why the stated hypothesis is structural rather than cosmetic.
\end{solution}

\begin{problem}[Algebraic stability proof dossier]\label{prob:iii14-dossier-10}
Prove the following structural statement and identify the step where its main hypothesis is used: If $f\in\mathcal S$, then $x^\alpha D^\beta f\in\mathcal S$.
\end{problem}
\begin{solution}
By the Leibniz rule, every Schwartz seminorm of this expression is bounded by a finite sum of higher Schwartz seminorms of $f$ with nonnegative constant coefficients. This proof also explains why the stated hypothesis is structural rather than cosmetic.
\end{solution}

\begin{problem}[Schwartz convolution closure proof dossier]\label{prob:iii14-dossier-11}
Prove the following structural statement and identify the step where its main hypothesis is used: If $f,g\in\mathcal S$, then $f*g\in\mathcal S$.
\end{problem}
\begin{solution}
Differentiate under the integral and distribute polynomial weights between $x-y$ and $y$; rapid decay of both factors controls every resulting seminorm. This proof also explains why the stated hypothesis is structural rather than cosmetic.
\end{solution}

\begin{problem}[Chapter synthesis]\label{prob:iii14-dossier-12}
Connect the main definitions and structural theorems of this chapter into one proof strategy. Explain what object is controlled, what estimate or closure principle is used, and what conclusion becomes available.
\end{problem}
\begin{solution}
The Schwartz space packages smoothness and rapid decay into a Fréchet-type test-function space stable under Fourier transformation. It is the natural common domain for classical transform calculus and tempered distributions. The recurring strategy is to encode the object in the chapter's natural measurable, normed, Fourier, or distributional structure, establish the controlling estimate, and then pass to the desired limit or representation.
\end{solution}
\section{Exercises with complete solutions}
The shorter exercise layer remains separate from the solved dossiers.

\begin{exercise}\label{exr:iii14-01}
State and apply the central principle behind Rapid decay. Give one immediate consequence in the setting of this chapter.
\end{exercise}
\begin{hint}
To prove Schwartz decay, bound \(x^\alpha\partial^\beta f(x)\) uniformly for every pair of multi-indices.
\end{hint}
\begin{solution}
A Schwartz function decays faster than every polynomial together with all derivatives. As an immediate consequence, the same principle can be inserted into later convergence, approximation, transform, or weak-form arguments whenever its hypotheses are verified.
\end{solution}

\begin{exercise}\label{exr:iii14-02}
State and apply the central principle behind Schwartz seminorms. Give one immediate consequence in the setting of this chapter.
\end{exercise}
\begin{hint}
Differentiation preserves the Schwartz class because it only changes the derivative index in the seminorms.
\end{hint}
\begin{solution}
Quantities $\sup_x|x^\alpha D^\beta f(x)|$ encode the topology. As an immediate consequence, the same principle can be inserted into later convergence, approximation, transform, or weak-form arguments whenever its hypotheses are verified.
\end{solution}

\begin{exercise}\label{exr:iii14-03}
State and apply the central principle behind Algebraic stability. Give one immediate consequence in the setting of this chapter.
\end{exercise}
\begin{hint}
Polynomial multiplication preserves Schwartz decay by shifting the polynomial weight index.
\end{hint}
\begin{solution}
Multiplication by polynomials and differentiation preserve Schwartz space. As an immediate consequence, the same principle can be inserted into later convergence, approximation, transform, or weak-form arguments whenever its hypotheses are verified.
\end{solution}

\begin{exercise}\label{exr:iii14-04}
State and apply the central principle behind Translation and modulation. Give one immediate consequence in the setting of this chapter.
\end{exercise}
\begin{hint}
For the Fourier transform, integrate by parts to trade powers of frequency for derivatives of the original function.
\end{hint}
\begin{solution}
Translations and frequency modulations preserve the class. As an immediate consequence, the same principle can be inserted into later convergence, approximation, transform, or weak-form arguments whenever its hypotheses are verified.
\end{solution}

\begin{exercise}\label{exr:iii14-05}
State and apply the central principle behind Fourier invariance. Give one immediate consequence in the setting of this chapter.
\end{exercise}
\begin{hint}
Smooth compactly supported functions are Schwartz already in physical space, and their Fourier transforms are Schwartz too. Compact support without smoothness does not imply Schwartz membership or rapid Fourier decay.
\end{hint}
\begin{solution}
The Fourier transform maps Schwartz space bijectively to itself. As an immediate consequence, the same principle can be inserted into later convergence, approximation, transform, or weak-form arguments whenever its hypotheses are verified.
\end{solution}

\begin{exercise}\label{exr:iii14-06}
State and apply the central principle behind Convolution. Give one immediate consequence in the setting of this chapter.
\end{exercise}
\begin{hint}
Use the Schwartz seminorms to quantify convergence instead of relying on pointwise or uniform convergence alone.
\end{hint}
\begin{solution}
Schwartz functions convolve to Schwartz functions. Convolution of a Schwartz function with an arbitrary L1 function is smooth but need not have rapid decay; all polynomial absolute moments of that L1 function are a sufficient additional condition for Schwartz decay. As an immediate consequence, the same principle can be inserted into later convergence, approximation, transform, or weak-form arguments whenever its hypotheses are verified.
\end{solution}

\begin{exercise}\label{exr:iii14-07}
State and apply the central principle behind Density. Give one immediate consequence in the setting of this chapter.
\end{exercise}
\begin{hint}
A function can be Schwartz without compact support; the Gaussian is the standard test case.
\end{hint}
\begin{solution}
Schwartz space is dense in $L^p$ for finite $p$ and in many Sobolev spaces. As an immediate consequence, the same principle can be inserted into later convergence, approximation, transform, or weak-form arguments whenever its hypotheses are verified.
\end{solution}

\begin{exercise}\label{exr:iii14-08}
State and apply the central principle behind Role in distribution theory. Give one immediate consequence in the setting of this chapter.
\end{exercise}
\begin{hint}
To show a map is continuous on Schwartz space, estimate each output seminorm by finitely many input seminorms.
\end{hint}
\begin{solution}
Its rapid decay makes pairing with polynomial-growth objects well-defined, leading to tempered distributions. As an immediate consequence, the same principle can be inserted into later convergence, approximation, transform, or weak-form arguments whenever its hypotheses are verified.
\end{solution}


% BEGIN VOL03-EXPANSION III14-exercises-01
\section{Graded supplementary exercises}
These exercises extend the preserved foundational set and are graded by mathematical role.

\subsection*{Standard computations}

\begin{exercise}[Gaussian]\label{exr:iii14-expand-01}
Is exp(-x\textasciicircum{}2) a Schwartz function?
\end{exercise}
\begin{hint}
Polynomially weighted derivatives are bounded.
\end{hint}
\begin{solution}
Yes.
\end{solution}

\begin{exercise}[Compact support]\label{exr:iii14-expand-02}
Is every smooth compactly supported function Schwartz?
\end{exercise}
\begin{hint}
Outside support it vanishes.
\end{hint}
\begin{solution}
Yes.
\end{solution}

\begin{exercise}[Polynomial]\label{exr:iii14-expand-03}
Is x\textasciicircum{}2 Schwartz?
\end{exercise}
\begin{hint}
Check rapid decay.
\end{hint}
\begin{solution}
No.
\end{solution}

\begin{exercise}[Derivative closure]\label{exr:iii14-expand-04}
If f is Schwartz, what about f'?
\end{exercise}
\begin{hint}
Use seminorms.
\end{hint}
\begin{solution}
It is Schwartz.
\end{solution}

\begin{exercise}[Polynomial multiplication]\label{exr:iii14-expand-05}
If f is Schwartz, what about x\textasciicircum{}3 f?
\end{exercise}
\begin{hint}
Shift polynomial weights.
\end{hint}
\begin{solution}
It is Schwartz.
\end{solution}

\subsection*{Proofs}

\begin{exercise}[Derivative stability]\label{exr:iii14-expand-06}
Prove differentiation preserves Schwartz space.
\end{exercise}
\begin{hint}
Each seminorm becomes another higher derivative seminorm.
\end{hint}
\begin{solution}
All remain finite.
\end{solution}

\begin{exercise}[Polynomial stability]\label{exr:iii14-expand-07}
Prove multiplication by a polynomial preserves Schwartz space.
\end{exercise}
\begin{hint}
Combine polynomial weights.
\end{hint}
\begin{solution}
Every new seminorm is controlled by finitely many original seminorms.
\end{solution}

\begin{exercise}[Fourier invariance]\label{exr:iii14-expand-08}
Outline the seminorm proof.
\end{exercise}
\begin{hint}
Convert frequency moments and derivatives to spatial derivatives and moments.
\end{hint}
\begin{solution}
All resulting expressions are integrable and have bounded transforms.
\end{solution}

\begin{exercise}[Convolution closure]\label{exr:iii14-expand-09}
Explain why convolution of two Schwartz functions is Schwartz.
\end{exercise}
\begin{hint}
Differentiate under the integral and distribute polynomial weights.
\end{hint}
\begin{solution}
Rapid decay of both factors controls every seminorm.
\end{solution}

\subsection*{Counterexamples and hypothesis testing}

\begin{exercise}[Rapid not compact]\label{exr:iii14-expand-10}
Give a Schwartz function without compact support.
\end{exercise}
\begin{hint}
Use a Gaussian.
\end{hint}
\begin{solution}
The Gaussian is nonzero everywhere but rapidly decaying.
\end{solution}

\begin{exercise}[L1 not Schwartz]\label{exr:iii14-expand-11}
Give an integrable function that is not Schwartz.
\end{exercise}
\begin{hint}
Use an interval indicator.
\end{hint}
\begin{solution}
$1_{[0,1]}$ is integrable on $\mathbb R$ but is not smooth, hence is not Schwartz.
\end{solution}

\begin{exercise}[Smooth not Schwartz]\label{exr:iii14-expand-12}
Give a smooth function that is not Schwartz.
\end{exercise}
\begin{hint}
Use the constant function 1.
\end{hint}
\begin{solution}
The constant function $1$ is smooth on $\mathbb R$, but $\sup_x|x\cdot1|=\infty$, so it is not Schwartz.
\end{solution}

\subsection*{Applications and investigations}

\begin{exercise}[Test space]\label{exr:iii14-expand-13}
Why is Schwartz space natural for Fourier distributions?
\end{exercise}
\begin{hint}
It is smooth, rapidly decaying, and Fourier-invariant.
\end{hint}
\begin{solution}
All key operations remain in one stable test space.
\end{solution}

\begin{exercise}[Regularization]\label{exr:iii14-expand-14}
Why approximate rough data by Schwartz functions?
\end{exercise}
\begin{hint}
They are dense in many Lp/Sobolev settings and easy to manipulate.
\end{hint}
\begin{solution}
Prove estimates for smooth approximants and pass to limits.
\end{solution}

\subsection*{Challenge problems}

\begin{exercise}[Translation]\label{exr:iii14-expand-15}
Why does translation preserve Schwartz space?
\end{exercise}
\begin{hint}
Expand powers of x using x-a plus a.
\end{hint}
\begin{solution}
Translated seminorms are controlled by finite sums of original seminorms.
\end{solution}

\begin{exercise}[Modulation]\label{exr:iii14-expand-16}
Why does multiplying by exp(2 pi i b x) preserve Schwartz space?
\end{exercise}
\begin{hint}
Differentiate with Leibniz; the modulation has unit magnitude.
\end{hint}
\begin{solution}
Each derivative is a finite sum of constants times derivatives of f, so weighted bounds remain finite.
\end{solution}
% END VOL03-EXPANSION III14-exercises-01
\section*{Chapter summary}
The chapter now supplies a canonical theorem-and-dossier layer for subsequent measure, Fourier, distributional, Sobolev, and PDE arguments.
